\subsection{商空间}

定义 3. 设 X 是拓扑空间，R 是 X 上的一个等价关系。X 按 R 的商空间是指集合 X/R 赋以 X 的拓扑按关系 R 的商所得的拓扑（§ 2，第 4 目，例 1）。

除非另有明确说明，每当我们把 X/R 作为拓扑空间谈论时，都应理解为指 X 按 R 的商空间。我们常把这个拓扑空间说成是把 X 中属于模 R 同一等价类的点等同起来所得到的空间。

设 \(\varphi\) 是典范映射 \(X \to X/R\)。由定义（§ 2，第 4 目，命题 6 及其推论），\(X/R\) 中的开（相应地，闭）集是使 \(\bar{\varphi}^{1}(A)\) 在 \(X\) 中为开（相应地，闭）集的集合 A；换言之，\(X/R\) 中的开（相应地，闭）集与 \(X\) 的关于 \(R\) 饱和的开（相应地，闭）子集一一对应，并且是这些子集的典范像。

命题 6. 设 X 是拓扑空间，R 是 X 上的一个等价关系，φ 是 X 到 X/R 上的典范映射；则 X/R 到拓扑空间 Y 中的一个映射 f 连续，当且仅当 f \(\circ\) φ 在 X 上连续。

这是 § 2 第 4 目命题 6 的一个特例；它表达了这样的事实：商拓扑是关于映射 \(\varphi\) 的终拓扑。

命题 6 表明，X/R 到 Y 中的连续映射与 X 到 Y 中的、在每个模 R 等价类上都取常值的连续映射之间有一一对应。

例。* 考虑等价关系 \(x \equiv y \pmod{1}\)（实直线 \(\mathbf{R}\) 上的等价关系）；\(\mathbf{R}\) 按这个关系的商空间称为一维环面，记作 \(\mathbf{T}\)。点 \(x \in \mathbf{R}\) 的等价类由所有的点 \(x + n\)（其中 \(n\) 取遍有理整数集 \(\mathbf{Z}\)）组成。由命题 6，\(\mathbf{T}\) 上的连续函数与 \(\mathbf{R}\) 上的以 1 为周期的连续函数之间有一一对应。我们将在第五章 § 1 中回到这个重要的例子。*

推论. 设 X、Y 是两个拓扑空间，R（相应地，S）是 X（相应地，Y）上的一个等价关系，且设 \(f: X \to Y\) 是与等价关系 R 和 S 相容的连续映射（《集合论》R，§ 5，第 8 目）；则映射 \(g: X/R \to Y/S\)（由 \(f\) 所诱导（《集合论》R，§ 5，第 8 目））是连续的。

这是商结构的一个一般性质的一个特例（《集合论》第四章，§ 2，第 6 目，判别准则 CST 20）。

命题 7（商空间的传递性）. 设 R 与 S 是拓扑空间 X 上的两个等价关系，使得 R 蕴含 S，且设 S/R 是商空间 X/R 上的商等价关系（《集合论》R，§ 5，第 9 目）。则典范双射 \((\mathbf{X}/\mathbf{R})/(\mathbf{S}/\mathbf{R}) \to \mathbf{X}/\mathbf{S}\) 是一个同胚。

这是终拓扑的传递性的一个特例（§ 2，第 4 目，命题 7。参见《集合论》第四章，§ 2，第 3 目，判别准则 CST 21）。

